\documentclass[11pt,a4paper]{article}
\usepackage[margin=2.4cm]{geometry}
\usepackage[T1]{fontenc}
\usepackage{booktabs}
\usepackage{array}
\usepackage{xcolor}
\usepackage{titlesec}
\usepackage{enumitem}
\usepackage{tcolorbox}
\usepackage[hidelinks]{hyperref}

\definecolor{good}{HTML}{1B7F3B}
\definecolor{warn}{HTML}{B26B00}
\definecolor{bad}{HTML}{B3261E}
\definecolor{grey}{HTML}{666666}
\definecolor{rule}{HTML}{2B4C7E}

\newcommand{\done}{\textcolor{good}{\textbf{DONE}}}
\newcommand{\prog}{\textcolor{warn}{\textbf{IN PROGRESS}}}
\newcommand{\todo}{\textcolor{grey}{\textbf{NOT STARTED}}}
\newcommand{\solid}{\textcolor{good}{\textbf{solid}}}
\newcommand{\weak}{\textcolor{warn}{\textbf{weak}}}
\newcommand{\unsup}{\textcolor{bad}{\textbf{unsupported}}}

\titleformat{\section}{\large\bfseries\color{rule}}{\thesection.}{0.6em}{}
\titleformat{\subsection}{\normalsize\bfseries}{}{0em}{}
\setlist[itemize]{leftmargin=1.1em,itemsep=2pt,topsep=3pt}
\setlength{\parskip}{5pt}
\setlength{\parindent}{0pt}
\renewcommand{\arraystretch}{1.25}

\newtcolorbox{phasebox}[1]{colback=white,colframe=rule!55,boxrule=0.7pt,
  arc=2pt,left=8pt,right=8pt,top=6pt,bottom=6pt,title={\bfseries #1},
  fonttitle=\color{white},colbacktitle=rule}

\title{\textbf{Plexus Discovery}\\[6pt]
\large Progress note --- written to be read, not decoded}
\author{}
\date{31 July 2026}

\begin{document}
\maketitle
\vspace{-2.2\baselineskip}

\begin{center}\small\color{grey}
Plain English throughout. Organised by phase, because phases are where we stop and you decide.\\
The technical specification is preserved separately in \texttt{plexus2\_technical.pdf}.
\end{center}

\vspace{4pt}
\hrule
\vspace{10pt}

%======================================================================
\section{What we are doing --- two things at once}

\textbf{Track A --- build the method.} An automated loop that searches for biological
\emph{mechanisms} rather than parameter values. It proposes a change to the model's structure
(``remove the pushing force'', ``add oriented cell division''), writes down what it expects
\emph{before} running, runs the simulation, measures, and records whether it was right. The
discipline that makes it worth building: a change of numbers can never masquerade as a new idea.

\textbf{Track B --- reproduce Okuda's result, as the proof that Track A works.} Okuda et al.\
(2018) grew a hollow ball of cells that spontaneously forms tubes, undulations and branches,
driven by a chemical pattern on its own surface. We are rebuilding that here. If the loop finds
the mechanism largely by itself, the method is demonstrated. \emph{The figure is not a side quest
--- it is the evidence.}

The two tracks feed each other: the reproduction is the honest test of the method, and the method
is what makes the reproduction more than hand-tuning.

%======================================================================
\section{Where we actually are}

Track A is further along than Track B, and until this morning I was over-reporting both.

An independent review this morning found that \textbf{three of the measurements the loop relies on
are broken}, and you have since found a fourth. Not the simulations --- the \emph{rulers}. A large part of what I reported over the
last two days was measured with faulty instruments and does not currently stand up.

\begin{tcolorbox}[colback=bad!4,colframe=bad!55,boxrule=0.7pt,arc=2pt]
\textbf{The clearest example, and why you were right to stop me.}

I ran 32 simulations overnight, sweeping five settings, and reported a strong conclusion:
\emph{``making the chemistry shape the tissue and keeping the tissue intact are mutually
exclusive.''}

Every one of those 32 runs ended with exactly \textbf{1778 cells}. I noticed that, wrote it down,
and called it ``remarkable''. I never asked the obvious question: \emph{why is it identical every
time?}

The answer is arithmetic. The memory buffer for the mesh holds 3552 points, and the geometry of a
closed cell sheet fixes the largest possible number of cells at $(3552+4)/2 = 1778$. \textbf{Every
run hit the wall and stopped.} I was not measuring biology. I was measuring the size of an array.

That conclusion is now marked \unsup. It may still turn out to be true --- but it is not measured,
and I presented it as if it were.
\end{tcolorbox}

%======================================================================
\section{What we believe, and how confident}

\begin{center}\small
\begin{tabular}{@{}p{7.6cm} l p{5.6cm}@{}}
\toprule
\textbf{Claim} & \textbf{Status} & \textbf{Why}\\
\midrule
The chemistry in the two front-page website movies does nothing to the shape
& \solid
& Arithmetic. The switch that lets chemistry drive growth is set to 50; the chemical only reaches
about 0.4. It never fires. The two movies are one simulation in two colour schemes.\\

Okuda's published settings are not reachable in our system today
& \solid
& Three of his values sit outside the ranges our search may use (\S\ref{sec:okuda}).\\

We can still regenerate the website movies from their recipes
& \weak
& They agree --- but old and new runs both stopped at the same 1778 ceiling, so the agreement
means less than it appeared to.\\

Our tube is ``forced'' rather than ``grown''
& \unsup
& The test for this was measuring the wrong thing entirely (\S\ref{sec:rulers}).\\

Chemistry-driven shaping destroys the cell sheet
& \unsup
& The 1778 problem. Needs re-running with a larger buffer --- that is Phase~1.\\
\bottomrule
\end{tabular}
\end{center}

%======================================================================
\section{The four broken rulers}\label{sec:rulers}

All four are \emph{measurement} faults. The physics may be fine; we simply could not see it.
The fourth is the camera, and you spotted it by looking at the movies --- it is the ruler the
automatic eye-check reads, and it is broken in two independent ways.

\begin{enumerate}[leftmargin=1.4em]
\item \textbf{The ``does the shape survive?'' test was measuring a fresh sphere.}
The intent: grow the tube, switch off growth and pushing, let it relax, see what survives. A real
tube survives; an artificially forced one collapses. The code was meant to continue from the end of
the run --- instead it started a brand-new simulation from the initial sphere, and so measured the
starting condition every time. It returned essentially the same value (1.014) in 14 of 16 runs.
\textbf{That value carries full weight in the score that ranks every experiment.}
\emph{Not yet fixed --- the one open item in Phase~0.}

\item \textbf{Two different quantities were both called ``how far it sticks out''.}
One measured distance from the centre \emph{of the tissue}; the other from the \emph{origin of the
coordinate system}. The ball drifts sideways as it grows, so the second reads drift as elongation.
\emph{Fixed this morning.}

\item \textbf{A safety check existed, but not where the danger was.}
We once chased a phantom result for days, caused by pairing mismatched frames. A guard was written
to catch it. My overnight study used exactly the pattern that guard forbids, in three places, and
never called the guard. \emph{Fixed this morning.}

\item \textbf{The camera hides the thing we are looking for, and hides growth entirely.}
\emph{(Your observation, confirmed in the code.)} Two separate faults, both in the movie every
automatic eye-check reads:
\begin{itemize}
\item \textbf{One fixed viewpoint.} The campaign's movies are drawn from a single angle that never
moves. A tube growing away from that angle sits \emph{behind the body} and is simply not in the
picture. There is no second view --- no top-down --- to catch it.
\item \textbf{The zoom follows the object.} The frame is re-fitted to the tissue on \emph{every
frame}, so a ball that doubles in size is drawn the same size throughout. \textbf{Growth is
invisible by construction.} Nothing that watches these movies --- human or machine --- can see the
thing the whole campaign is about.
\end{itemize}
Okuda's figures use a fixed frame and a consistent viewpoint, which is exactly why they read
clearly. \emph{Not fixed --- now part of Phase~0, because it is an instrument, not a nicety.}
\end{enumerate}

%======================================================================
\section{What is genuinely working}

The corrections should not hide the parts that earned their keep.

\begin{itemize}
\item \textbf{Predictions are written down before runs, and wrong ones are recorded as wrong.}
Twice in two days I proposed a mechanism and the measurement killed it within the hour. Both sit in
the record as retractions rather than being quietly edited away. This is the part I would defend
hardest.
\item \textbf{The video check found a disagreement --- but I over-claimed who was right.}
On the first full round the system described each movie in words automatically and compared that
to the numbers. It \emph{vetoed the two highest-scoring runs}, including one that three separate
analyses had called a tube. I reported this as ``the numbers were wrong and the picture was
right''. \textbf{That was too strong.} Given ruler~4 above, the picture may have been hiding a
tube behind the body, or failing to show growth at all. What the check genuinely established is
that \emph{the two instruments disagree} --- which is still valuable, and is precisely why it
exists. Which one to believe is not yet decided, and cannot be until the camera is fixed.
\item \textbf{The loop refuses to guess.} Asked to score 32 results against predictions written in
a unit it did not recognise, it declined 32 times rather than inventing 32 confirmations.
\end{itemize}

%======================================================================
\section{The team of agents --- what is actually switched on}\label{sec:agents}

You were told 6 of 12 are implemented or used. I audited it directly against the code rather than
accept the figure, and it is better than that --- but there are two real gaps, and they are worth
closing exactly as you say.

\textbf{The design calls for 15 roles} (14 automatic, plus a human who patches the substrate from
outside the loop). \textbf{All 15 are written. 13 actually run.} The two that do not:

\begin{center}\small
\begin{tabular}{@{}l p{6.2cm} p{6.4cm}@{}}
\toprule
\textbf{Role} & \textbf{What it is for} & \textbf{Status}\\
\midrule
\textbf{Evolution}
& Take the round's winner and improve it --- simplify it, combine it with a runner-up, ground it
in the paper. Without it the loop can only ever \emph{pick} from what it happened to propose; it
can never \emph{refine}.
& Written, \textcolor{bad}{\textbf{never called anywhere}}.\\
\textbf{Referee}
& Rank a batch by a proper tournament (every candidate against every other, aggregated) rather
than by sorting on one number. The tournament is what protects against ordering bias.
& Written, but \textcolor{bad}{\textbf{only ever exercised in its own self-test}}. The live round
sorts on a single score instead.\\
\bottomrule
\end{tabular}
\end{center}

The other thirteen --- the one that holds the objective, the one that reads the papers, the one
that proposes, the peer-reviewer, the type-checker, the second-opinion judge, the duplicate
detector, the three independent analysts, the eye-check, the one that writes the causal story, the
pattern-distiller, and the one that certifies the instruments --- all ran on the last round.

Two honest caveats: the eye-check is currently reading a broken camera (ruler~4), and the analysts'
budget ceiling is not enforced, so a round can overrun its time allowance. Both are on the list.

%======================================================================
\section{What stands between us and the Okuda figure}\label{sec:okuda}

Comparing his published settings against what our system can express: \textbf{three of his values
are outside the range our search is allowed to use}, so a faithful reproduction is impossible today
however long we run.

\begin{center}\small
\begin{tabular}{@{}l l l l@{}}
\toprule
\textbf{Okuda's setting} & \textbf{His value} & \textbf{Ours allows} & \textbf{Problem}\\
\midrule
Growth-switch sharpness       & 10              & 1--8    & out of range\\
How far the inhibitor spreads & 10              & 0.1--2  & out of range\\
Pattern length-scale          & 0.001--0.1      & 1--10   & \textcolor{bad}{not even the same quantity}\\
Chemistry speed               & 4 powers of ten & 1.2     & far too narrow\\
\bottomrule
\end{tabular}
\end{center}

There is also a subtler gap. In Okuda's model the chemistry responds to the tissue's \emph{shape}
--- chemical flows between neighbouring cells in proportion to how much wall they share. Ours
ignores geometry entirely. \emph{One half of the feedback loop is simply missing.}

\textbf{Decision already taken:} calibrate to Okuda's \emph{observations} --- how many spots
appear, how thick the tube is --- rather than copying his parameter values, which live in a
differently-scaled model.

%======================================================================
\section{The phases}

Eight stages. \textbf{I stop at the end of each one and wait for you.} That is the partition.

\begin{phasebox}{Phase 0 --- Fix the broken rulers \hfill \prog\ (2 of 4)}
\textbf{Goal.} Make the measurements trustworthy. Nothing runs unattended until this lands.

\textbf{Why.} Running before this writes uninterpretable results into a permanent record. Four
rulers were wrong; two are now right.

\textbf{Remaining.}
\begin{itemize}
\item \textbf{The ``does the shape survive?'' test} (\S\ref{sec:rulers}, item~1). Either repair it
so it continues from the end state, or remove it from the score. A constant carrying full weight is
worse than no term at all.
\item \textbf{The camera} (\S\ref{sec:rulers}, item~4). Fix the frame so it does not follow the
object --- one size, held for the whole run and shared across runs, as in Okuda's figures --- and
add a second viewpoint so a tube cannot hide behind the body. This was originally scheduled much
later as a presentation nicety. It is not: it is the ruler the eye-check reads, and until it is
fixed we cannot say whether the eye-check's vetoes were right.
\end{itemize}

\textbf{What you get.} A one-line confirmation that the rulers are sound, plus one before/after
movie so you can see the camera change with your own eyes.
\end{phasebox}

\begin{phasebox}{Phase 1 --- Re-measure what Phase 0 invalidated \hfill \todo\ ($\sim$1 hour)}
\textbf{Goal.} Re-run the 32-simulation study with a buffer ten times larger and the frame-pairing
guard live.

\textbf{Why.} This is the phase that decides whether my overnight conclusion was real.

\textbf{What it will decide.} One of two outcomes, and \emph{both are good}:
\begin{itemize}
\item it survives --- we have a genuine structural result, which is exactly what a paper about the
method wants; or
\item it dissolves --- and we avoided publishing an artefact.
\end{itemize}

\textbf{What you get.} The answer in a sentence, plus the corrected table.
\end{phasebox}

\begin{phasebox}{Phase 2 --- Make Okuda's settings reachable \hfill \todo\ ($\sim$1 week)}
\textbf{Goal.} Widen the three out-of-range settings; replace the mislabelled length-scale with the
quantity Okuda actually varies; add the missing shape-to-chemistry feedback; add a fixture that
asserts his published constants rather than tuning them.

\textbf{Why.} Without this, Track B cannot succeed at any effort level.

\textbf{What you get.} A list of what now matches his paper and what still does not, with reasons.
\end{phasebox}

\begin{phasebox}{Phase 2b --- Switch on the two agents that never run \hfill \todo\ (1 day)}
\textbf{Goal.} Put \textbf{Evolution} and the \textbf{Referee} into the live round
(\S\ref{sec:agents}), and enforce the analysts' time budget.

\textbf{Why.} Without Evolution the loop can only pick among the changes it happened to think of,
never improve the best one --- which is half of what makes it a search rather than a lottery.
Without the Referee's tournament, ranking rests on a single number, and that number is the one
Phase~0 is currently repairing.

\textbf{What you get.} Confirmation that all fourteen automatic roles ran in one round, with the
round's cost in time and tokens.
\end{phasebox}

\begin{phasebox}{Phase 3 --- Measurements that can tell his regimes apart \hfill \todo\ (2 days)}
\textbf{Goal.} Four measurements: how many spots, how thick the tube, how many branches over time,
and the shape of the surface. Calibrate against his roughly five spots at 2000 cells, then freeze
the calibration.

\textbf{Why.} We currently have no measurement that can distinguish a tube from an undulation from
a branch --- precisely the three outcomes the figure has to show.

\textbf{What you get.} The four numbers, checked against his published figures.
\end{phasebox}

\begin{phasebox}{Phase 3b --- Point the camera \emph{along the tube} \hfill \todo\ (half a day)}
\textbf{Goal.} Phase~0 makes the camera honest --- fixed frame, second viewpoint. This phase makes
it \emph{smart}: lock it to the tube's own axis, one view side-on and one looking straight down the
tube.

\textbf{Why.} The end-on view \emph{is} the thickness and branch-count measurement, so camera and
measurement share one axis. It is a refinement, not a repair --- the repair is in Phase~0.

\textbf{What you get.} Movies in which the tube is always in full profile, whichever way it grew.
\end{phasebox}

\begin{phasebox}{Phase 4 --- The demonstration \hfill \todo\ (one wave, $\sim$1 hour)}
\textbf{Goal.} Run the designed grid --- six to eight simulations --- that produces the figure:
tube, undulation and branching, at his settings.

\textbf{What you get.} \emph{The figure.}
\end{phasebox}

\begin{phasebox}{Phase 5 --- The method claim \hfill \todo}
\textbf{Goal.} Write up what the loop found by itself versus what we told it, with the record of
predictions and retractions as the evidence.

\textbf{What you get.} The argument that Track A works, resting on Track B.
\end{phasebox}

%======================================================================
\section{What I need from you}

Nothing blocking. Two things worth knowing:

\begin{enumerate}[leftmargin=1.4em]
\item \textbf{Phase 1 is the interesting one} and it costs about an hour of compute. When it
finishes I will bring you the answer, not the reasoning --- unless you want the reasoning, in which
case say so and it becomes the default.
\item \textbf{Tell me if this note is the right shape.} Too long, too short, wrong order --- it is
cheap to change, and it is the thing that lets you steer.
\end{enumerate}

%======================================================================
\section*{Appendix --- words I could not avoid}

Only these four. Everything else above is ordinary English.

\begin{itemize}
\item \textbf{Cell sheet (or mesh)} --- the model of the tissue: a hollow ball of polygonal cells
sharing walls, like a football.
\item \textbf{Operator} --- one mechanism as a reusable building block: ``cells divide'',
``chemical spreads'', ``surface pulls tight''. A model is a combination of these.
\item \textbf{The loop} --- the automated cycle: propose a change, predict the outcome, run,
measure, record whether the prediction held.
\item \textbf{Phase} --- one of the eight stages above. Our agreed stopping points.
\end{itemize}

\vspace{6pt}
\hrule
\vspace{6pt}
{\small\color{grey}
Technical detail --- the control law, the agent roles, the operator language, the validation
ladder and the full external review --- is preserved in \texttt{plexus2\_technical.pdf} (23 pp).
My working log is \texttt{SESSION\_LOG.md}. You should not need either.}

\end{document}
